\documentclass[a4paper]{article}

\usepackage{graphicx}
\usepackage{amsmath,amssymb}
\usepackage{minted}
\usepackage{multirow}
\usepackage{float}
\usepackage{algorithm}
\usepackage{algpseudocode}
\usepackage{todonotes}
\usepackage{forest}
\usepackage{xstring}
\usepackage{xcolor}
\usepackage[most]{tcolorbox}
\usepackage{xparse}
\usepackage{natbib}
\usepackage{adjustbox}

\usetikzlibrary{graphs,graphdrawing}
\usegdlibrary{layered}

% for external graphviz
\input{/home/donkarlo/Dropbox/repo/nd_language_ai_project/src/nd_language_ai/formal/computer/programming/tex/latex/code_clips/external_graphviz.tex}

% for tags
\input{/home/donkarlo/Dropbox/repo/nd_language_ai_project/src/nd_language_ai/formal/computer/programming/tex/latex/part/control_sequence/kind/semtag/semtag.tex}

% for links
\input{/home/donkarlo/Dropbox/repo/nd_language_ai_project/src/nd_language_ai/formal/computer/programming/tex/latex/code_clips/seehere.tex}

% for nested lists and sections
\input{/home/donkarlo/Dropbox/repo/nd_language_ai_project/src/nd_language_ai/formal/computer/programming/tex/latex/code_clips/deep_structure.tex}

\title{Isolation Forest}
\date{}
\author{Mohammad Rahmani}

\begin{document}
   \maketitle

\section{Scoring rule}
Each isolation tree repeatedly selects a feature and a split uniformly within that feature's observed range in the current node. Trees are built on subsamples. A point that reaches a leaf after few splits tends to be easier to isolate. Average path lengths across the forest, using the expected extra path length when a terminal node still contains multiple training samples.

\section{Example}
In a dataset of transaction amounts and counts, a transaction with a much higher amount than most others can be isolated by a few splits on amount. Transactions near the center of a dense group tend to require more splits.

\section{Practical considerations}
Tune the number of trees for stable rankings and select a threshold on validation data or from a review budget. Check the score direction in the implementation. Many irrelevant features can dilute the signal, and axis-aligned splits can struggle with anomalies defined by oblique relationships.



   \bibliographystyle{plainnat}
   \bibliography{/home/donkarlo/Dropbox/repo/nd_shared_research/bibliography/bibliography.bib}
\end{document}
